\documentclass[11pt, a4paper]{article}

% Helpful packages
\usepackage{geometry}
\usepackage{indentfirst}
\usepackage{enumitem}
\usepackage[cjk]{kotex}
\usepackage{hyperref}
\usepackage{setspace}
\usepackage[T1]{fontenc}
\usepackage{relsize}
\usepackage{anyfontsize}
\usepackage{xcolor}

\fontfamily{georgia}\selectfont

% Global formatting tweaks
\pagestyle{empty}
\geometry{margin=0.5in}
\setlength\parindent{0pt}
\setlist{nosep}
\renewcommand\labelitemi{--}

% New commands
\newcommand{\name}[1]{
    \begin{center}
    \section*{\huge #1}
    \end{center}
}

\newcommand{\topinfo}[1]{
    \begin{center}
    \vspace{-0.2cm}#1
    \vspace{-0.1cm}
    \end{center}
}

\newcommand{\resumesection}[1]{
    \vspace{-0.35cm}
    \section*{#1}
    \vspace{-0.25cm}
    \hrule\vspace{0.2cm}
}

\hypersetup{
    colorlinks=true,
    linkcolor=blue,
    urlcolor=blue,
    pdfborder={0 0 1},
}

\definecolor{darknavy}{RGB}{0,45,140}
\definecolor{deepnavy}{RGB}{0,25,100}


\begin{document}

\name{Jihwan Oh}

\onehalfspacing
\topinfo{Stanford University, Stanford, CA, USA}

\onehalfspacing
\topinfo{
\href{mailto:jihwanoh@stanford.edu}{jihwanoh@stanford.edu}
\; | \;
\href{https://www.linkedin.com/in/jihwan-oh-21a772305}{
\underline{\textcolor{blue}{LinkedIn}}}
\; | \;
\href{https://jihwan01.github.io/}{
\underline{\textcolor{blue}{jihwan01.github.io}} (Website)}
\; | \;
\href{https://github.com/jihwan01}{
\underline{\textcolor{blue}{GitHub}}}
}

\setstretch{1.178}


%%%%%%%%%%%%%%%%%%%%%%%%%%%%%%%%%%%%%%%%%%%%%%%
\resumesection{RESEARCH INTERESTS}

\begin{itemize}[leftmargin=1.2em]
    \item Cross-layer optimization and hardware-software co-design for energy-efficient, high-performance computing
    \item Computer Architecture, Specialized Accelerators, and Systems for ML
\end{itemize}


%%%%%%%%%%%%%%%%%%%%%%%%%%%%%%%%%%%%%%%%%%%%%%%
\resumesection{EDUCATION}

\begin{itemize}[leftmargin=1.2em]
    \item \textbf{Stanford University}
    \hfill Sep.\ 2026 --\\[-0.2em]
    {\smaller[0.6] Ph.D. Student in Electrical Engineering; Advisor: Prof.\ Thierry Tambe}
\end{itemize}

\begin{itemize}[leftmargin=1.2em]
    \item \textbf{Korea Advanced Institute of Science and Technology (KAIST)}
    \hfill Feb.\ 2019 -- Feb.\ 2026\\[-0.2em]
    {\smaller[0.6] B.S. in School of Electrical Engineering (Primary Major), School of Computing (Double Major)}\\[-0.2em]
    {\smaller[0.6] GPA: 4.07/4.3, \textit{Summa Cum Laude}; Dean's List: Spring 2019, Spring 2024}
\end{itemize}

\begin{itemize}[leftmargin=1.2em]
    \item \textbf{Georgia Institute of Technology}
    \hfill Jan.\ 2025 -- Jul.\ 2025\\[-0.2em]
    {\smaller[0.6] Exchange Program, School of Electrical and Computer Engineering; GPA: 4.0/4.0}
\end{itemize}


%%%%%%%%%%%%%%%%%%%%%%%%%%%%%%%%%%%%%%%%%%%%%%%
\resumesection{PUBLICATIONS}

% ----- IISWC 2026 -----
\noindent
\begin{minipage}[t]{0.85\textwidth}
\textbf{Compute-Communication Overlap Is Not Free: A Cross-Layer\\
Characterization in GPU LLM Workloads}
\end{minipage}%
\begin{minipage}[t]{0.15\textwidth}
\raggedleft Sep.\ 2026
\end{minipage}\\[-0.3em]

{\smaller[0.6]
\underline{\textbf{\textit{Jihwan Oh}}},
Seokjin Go, Junkyum Kim, Jongse Park, Divya Mahajan
}

{\smaller[0.6]
\textit{Accepted to IEEE International Symposium on Workload Characterization (IISWC) 2026}; \textbf{Best Paper Finalist}
}

\vspace{0.1cm}


% ----- ISPASS 2025 -----
\noindent
\parbox[t]{0.88\textwidth}{%
\textbf{Characterizing Compute-Communication Overlap in GPU-Accelerated\\
Distributed Deep Learning: Performance and Power Implications}
}%
\hfill
\parbox[t]{0.10\textwidth}{%
\raggedleft May.\ 2025
}\\[-0.3em]

{\smaller[0.6]
Seonho Lee,
\underline{\textbf{\textit{Jihwan Oh}}},
Junkyum Kim, Seokjin Go, Jongse Park, Divya Mahajan
}

{\smaller[0.6]
\textit{Poster presented at ISPASS 2025},
\href{https://arxiv.org/abs/2507.03114}{arXiv:2507.03114}
}


%%%%%%%%%%%%%%%%%%%%%%%%%%%%%%%%%%%%%%%%%%%%%%%
\resumesection{RESEARCH EXPERIENCE}

\textbf{Systems Infrastructure and Architecture Research Lab, Georgia Tech}
\hfill Jan.\ 2025 -- Aug.\ 2026

\begin{description}[wide=0pt]
    \item \textit{Researcher (Advisor: Prof.\ Divya Mahajan)}
\end{description}

\begin{itemize}[leftmargin=2em]

    \item Characterized compute-communication overlap in large-scale distributed LLM training and showed that the widely held assumption ``maximizing overlap is always beneficial'' breaks down, with aggressive overlap degrading both execution time and power efficiency across parallelism configurations. Published and presented as a poster at ISPASS 2025.

    \item Led a project to identify the hardware-level causes of compute-communication overlap overheads that remained hidden at the software layer through profiling and analysis.

    \item Designed a new NCCL cross-layer communication protocol that repurposes unused shared memory as a TMA-driven staging buffer to enable internal pipelining, achieving comparable or higher bandwidth with roughly 50\% fewer SMs in microbenchmarks.

\end{itemize}

\vspace{0.18cm}

\textbf{Computer Architecture and Systems Lab, KAIST}
\hfill Sep.\ 2024 -- Dec.\ 2024

\begin{description}[wide=0pt]
    \item \textit{Undergraduate Researcher (Advisor: Prof.\ Jongse Park)}
\end{description}

\begin{itemize}[leftmargin=2em]

    \item Analyzed LLM inference on \textit{NeuPIMs} (ASPLOS 2024) and found that FFN compute grows quadratically with hidden size, outpacing other operations and causing severe NPU-PIM load imbalance that limits model-size scalability. Designed a redistribution strategy that balances FFN execution across NPU and PIM, improving utilization and end-to-end performance.

    \item Supported research on applying quantization techniques to PIM architectures by surveying state-of-the-art methods and benchmarking accuracy across multiple models.

\end{itemize}

\vspace{0.18cm}

\textbf{Urban Robotics Lab, KAIST}
\hfill Feb.\ 2024 -- Jun.\ 2024

\begin{description}[wide=0pt]
    \item \textit{Undergraduate Researcher (Advisor: Prof.\ Hyun Myung)}
\end{description}

\begin{itemize}[leftmargin=2em]
    \item Implemented SLAM algorithms and ROS-based motion planning systems for mobile robotics.
\end{itemize}


%%%%%%%%%%%%%%%%%%%%%%%%%%%%%%%%%%%%%%%%%%%%%%%
\resumesection{TEACHING \& INDUSTRY EXPERIENCE}

\textbf{CS.10001: Introduction to Programming, KAIST}, Teaching Assistant
\hfill Sep.\ 2025 -- Dec.\ 2025

\vspace{0.05cm}

\textbf{Yulmae Academy}, Programming Instructor
\hfill Dec.\ 2023 -- Dec.\ 2024

\begin{itemize}[leftmargin=1.2em]
    \item Taught C, Python, and introductory machine learning to science high school students.
\end{itemize}

\vspace{0.05cm}

\textbf{CS.93000: Immersion Camp (Coding Camp), KAIST}, Teaching Assistant
\hfill Oct.\ 2023 -- Feb.\ 2024

\vspace{0.05cm}

\textbf{Republic of Korea Air Force}, Software Developer
\hfill Aug.\ 2021 -- May.\ 2023

\begin{itemize}[leftmargin=1.2em]

    \item Developed and refined a VR flight-training system using Unreal Engine 4 and C++, integrating pilot feedback and Git-based collaboration to enhance realism and training effectiveness.

    \item Selected to present the simulator as the Air Force representative at the 2022 ROKAF Information \& Communication Development Conference.

    \item Selected to serve as a counseling officer supporting soldier adaptation, dormitory operations, and mental-health guidance.

\end{itemize}


%%%%%%%%%%%%%%%%%%%%%%%%%%%%%%%%%%%%%%%%%%%%%%%
\resumesection{ACADEMIC ACTIVITIES}

\textbf{Paper Presenter}, IISWC 2026 -- Boulder, CO, USA
\hfill Sep.\ 2026

\begin{itemize}[leftmargin=1.2em]
    \item Presenting the first-author paper ``Compute-Communication Overlap Is Not Free: A Cross-Layer Characterization in GPU LLM Workloads''.
\end{itemize}

\textbf{Poster Presentation}, ISPASS 2025 -- Ghent, Belgium
\hfill May.\ 2025

\begin{itemize}[leftmargin=1.2em]
    \item Presented poster on ``Characterizing Compute-Communication Overlap in GPU-Accelerated Distributed Deep Learning''.
\end{itemize}

\textbf{Attendee}, ISCA 2025 and uArch Workshop -- Tokyo, Japan
\hfill Jun.\ 2025

\begin{itemize}[leftmargin=1.2em]
    \item Selected as a full travel grant recipient to attend ISCA 2025 and participate in the uArch Workshop.
\end{itemize}


%%%%%%%%%%%%%%%%%%%%%%%%%%%%%%%%%%%%%%%%%%%%%%%
\resumesection{AWARDS \& SCHOLARSHIPS}

\begin{itemize}[leftmargin=1.2em]

    \item \textbf{Stanford Thrive Fellow}, Stanford University
    \hfill 2026

    \vspace{0.03cm}

    \item \textbf{Best Paper Finalist}, IEEE IISWC 2026
    \hfill 2026

    \vspace{0.03cm}

    \item \textbf{Next-Generation Engineer Award: Highest Distinction}
    \hfill 2025\\
    Institute for Promotion of Engineering and Science of Korea, IPESK\\
    Awarded the top distinction as the single highest-ranked recipient in the department.

    \vspace{0.03cm}

    \item \textbf{uArch Mentoring Workshop Full Travel Grant}, uArch @ ISCA 2025
    \hfill 2025\\
    Awarded a fully funded travel grant for both the uArch workshop and ISCA 2025.

    \vspace{0.03cm}

    \item \textbf{Student Travel Grant}, IEEE ISPASS 2025
    \hfill 2025

    \vspace{0.05cm}

    \item \textbf{Dean's List}, School of Electrical Engineering, KAIST
    \hfill 2024

    \vspace{0.03cm}

    \item \textbf{Korea-U.S. Student Exchange Program Scholarship}
    \hfill 2024\\
    Korea Institute for Advancement of Technology, KIAT

    \vspace{0.05cm}

    \item \textbf{National Science and Technology Scholarship}, Korea Government
    \hfill 2021 -- 2024

    \vspace{0.03cm}

    \item \textbf{Dean's List}, Freshman Division, KAIST
    \hfill 2019

    \vspace{0.05cm}

    \item \textbf{Minister of Education Award}, Ministry of Education (STEAM R\&E Program)
    \hfill 2018

\end{itemize}


%%%%%%%%%%%%%%%%%%%%%%%%%%%%%%%%%%%%%%%%%%%%%%%
\resumesection{TECHNICAL SKILLS}

\begin{itemize}[leftmargin=1.2em]

    \item \textbf{Programming Languages:}
    C/C++, Python, Rust, SystemVerilog, Verilog

    \vspace{0.05cm}

    \item \textbf{Libraries and Tools:}
    CUDA, PyTorch, NCCL, CUTLASS, cuBLAS;
    Megatron-LM, Megatron-DeepSpeed;
    Nsight Compute, Nsight Systems;
    ROS, Unreal Engine 4; Git

\end{itemize}

\end{document}
